\subsubsection{Higher-Order Judgments in Twelf}
Inference rules involving typing contexts
such as rule \code{T.6} are \emph{hypothetical judgments}~\cite{TwelfHypothetical}
or judgments that make use of hypothetical assumptions.
Hypothetical judgments can be encoded in Twelf using
\emph{higher-order judgments}~\cite{TwelfHigherOrder}
or judgments that contain other judgments.
Higher-order judgments are encoded in Twelf as \emph{higher-order types}.
Higher-order types are functions types, where one of their input types
is also a function type.
Encoding higher-order judgments as higher-order types is similar to
the technique of higher-order abstract syntax discussed in
Section~\ref{sec:twelf-hoas}, where syntactic terms with binders
(e.g. the body of the \code{let} expression) are encoded as
terms with holes or equivalently
functions that take in terms as input and return terms.
Higher-order judgments are encoded with function types that have
input types representing the hypothetical assumptions.

Higher-order types representing hypothetical judgments involve
\emph{pi-types}, denoted $\Pi x:S.T$.
A pi-type is a special kind of function type.
Terms of pi-types are \emph{pi-abstractions}.
First, a \emph{lambda-abstraction}, denoted ``$\lambda x:S.e$'',
is a function term of a more conventional kind of function type,
 $\ftype{S}{T}$.
A function of this type takes in a term of type $S$ and
returns a term of type $T$.
However, a pi-abstraction of pi-type $\Pi x:S.T$ would map
a term $s$ of type $S$ to a term of type $[s / x]T$.
That is, the return type of a pi-type can vary according
to the argument supplied.
Hence, if $x$ is (syntactically) a part of $T$ in the pi-type
$\Pi x:S.T$, then a term returned by a function of that
pi-type is a term of a dependent type, since the return type
$T$ depends on the input term $x$.
If $x$ is not a part of $T$ in $\Pi x:S.T$,
then we abbreviate $\Pi x:S.T$ as $\ftype{S}{T}$,
which also signals that a lambda-abstraction can be of
this type, since the return type does not involve the argument.
Lastly, the pi-type $\Pi x:S.T$ is represented in Twelf syntax
as \code{\{x:S\} T}.

A common Twelf coding convention is used in
the return type of the \code{of/let} function term
given on line 20 of the \code{typing.elf} file:
``\mbox{\code{of (let E1 ([x] E2 x)) T2}}''.
This code snippet could have been replaced with
the shorter snippet:
``\mbox{\code{of (let E1 E2) T2}}''.
The \code{E2} in the shorter snippet
already denotes a term
of function type ``\mbox{\code{exp -> exp}}''
(or equivalently, an \code{exp} with free variable
occurrences).
However, we applied a Twelf coding convention when
we replaced \code{E2} with the equivalent function
term \code{([x] E2 x)}.
The wrapper function \code{([x] E2 x)} is used 
just to make it easier to see that \code{E2} is
a function term that takes in a single argument.

The \code{of/let} term represents rule \code{T.6}
and is of a higher-order type:
\begin{verbatim}
of/let :  ({x: exp} of x T1 -> of (E2 x) T2) ->
          of E1 T1 ->
          of (let E1 ([x] E2 x)) T2
\end{verbatim}
The second argument type ``\mbox{\code{of E1 T1}}''
corresponds to the premise ``$e_1 : \tau_1$''.
The type of \code{of/let}'s first argument is the pi-type:
``\mbox{\code{\{x: exp\} of x T1 -> of (E2 x) T2}}'';
let $P$ denote this pi-type.
A pi-abstraction of type $P$ takes in two arguments:
The first is an \code{exp} term, which will be bound to
\code{x}.
The second is a term of type ``\mbox{\code{of x T1}}''
representing a proof that \code{x} has type \code{T1}.
Given two such terms,
a pi-abstraction of type $P$ should return
a derivation of ``\mbox{\code{of (E2 x) T2}}'',
where \code{(E2 x)} is the \code{E2} term with its hole
filled in by the \code{exp} term that is bound to \code{x}.

The pi-type $P$ represents the second premise
``\mbox{$\Gamma , x : \tau_1 \vdash e_2 : \tau_2$}''
of rule ~\code{T.6}.
Passing a derivation \code{dx} of type ``\mbox{\code{of x T1}}''
to a pi-abstraction simulates extending the typing context
with the assumption ``\mbox{$x : \tau_1$}''.
The derivation \code{dx} of type ``\mbox{\code{of x T1}}''
can be used in the body of the pi-abstraction to
return a term/proof of ``\mbox{\code{of (E2 x) T2}}''.
